\documentclass[12pt,a4paper]{article}

% -------------------- Packages --------------------
\usepackage[margin=1in]{geometry}
\usepackage{amsmath}
\usepackage{setspace}
\usepackage{times}
\usepackage{titlesec}
\usepackage{booktabs}
\usepackage{array}
\usepackage{longtable}
\usepackage{hyperref}
\usepackage{xcolor}
\usepackage{listings}

% -------------------- Formatting --------------------
\onehalfspacing
\setlength{\parindent}{0pt}
\setlength{\parskip}{6pt}

\titleformat{\section}
  {\normalfont\Large\bfseries}
  {\thesection.}{0.5em}{}

\titleformat{\subsection}
  {\normalfont\large\bfseries}
  {\thesubsection.}{0.5em}{}

\hypersetup{
    colorlinks=true,
    linkcolor=black,
    urlcolor=blue,
    citecolor=black
}

\lstdefinestyle{jsonstyle}{
    basicstyle=\ttfamily\small,
    breaklines=true,
    frame=single,
    backgroundcolor=\color{gray!5},
    showstringspaces=false
}

\lstset{style=jsonstyle}

% -------------------- Document --------------------
\begin{document}

% -------------------- Title Page --------------------
\begin{titlepage}
    \centering

    \vspace*{1.5cm}

    {\Large \textbf{Galala University}}\\[1.5cm]

    {\LARGE \textbf{Civara}}\\[0.4cm]

    {\Large \textbf{Frontend Application Requirements}}\\[0.5cm]

    {\large \textbf{A Web Dashboard for Student Performance Statistical Analysis}}\\[2cm]

    \begin{table}[h!]
        \centering
        \renewcommand{\arraystretch}{1.4}
        \begin{tabular}{>{\bfseries}l l}
            Project Name: & Civara \\
            Domain: & \texttt{civara.mahros.dev} \\
            API Base (Prod): & \texttt{https://civara.mahros.dev/api/v1} \\
            API Base (Local): & \texttt{http://localhost:3030/api/v1} \\
        \end{tabular}
    \end{table}

    \vfill

    {\large \today}
\end{titlepage}

% -------------------- Main Content --------------------

\section{Overview}

This document defines the requirements for a frontend web application compatible with the Civara API contract (\texttt{requirements/api-contract.tex}) and the Postman collection (\texttt{postman.json}).

The frontend provides a simple, professional dashboard that allows users to:
\begin{itemize}
    \item Upload a CSV dataset of student scores.
    \item Choose the score type to analyze.
    \item Submit the dataset to the API for analysis.
    \item View and interpret results (descriptive statistics, confidence interval, hypothesis test summary).
    \item Cache the latest analysis locally to support fast reloads without a database.
\end{itemize}

\section{Non-Goals}

\begin{itemize}
    \item No authentication, login, roles, or user accounts.
    \item No server-side persistence requirements (the API does not store datasets).
    \item No admin panel.
    \item No multi-user collaboration.
\end{itemize}

\section{Target Users}

\begin{itemize}
    \item Students learning descriptive statistics and hypothesis testing.
    \item Instructors demonstrating statistical inference with real datasets.
    \item General users exploring the student performance dataset structure.
\end{itemize}

\section{Supported Browsers and Devices}

\begin{itemize}
    \item Desktop: Latest Chrome, Edge, and Firefox.
    \item Mobile: Responsive layout is required, but data tables may switch to stacked cards for readability.
\end{itemize}

\section{Frontend Architecture (High-Level)}

\subsection{Recommended Technology}

The implementation must be a modern single-page application (SPA). The exact framework may vary, but the minimum expectations are:
\begin{itemize}
    \item TypeScript support (recommended).
    \item Component-based UI architecture.
    \item A production build pipeline with environment-based configuration.
\end{itemize}

\subsection{Configuration}

The frontend must use a configurable API base URL.

Required configuration variable:
\begin{itemize}
    \item \texttt{API\_BASE\_URL} (e.g., \texttt{https://civara.mahros.dev/api/v1}).
\end{itemize}

For local development:
\begin{itemize}
    \item \texttt{API\_BASE\_URL = http://localhost:3030/api/v1}
\end{itemize}

\section{API Integration Requirements}

\subsection{Endpoints}

The frontend must integrate with the following endpoints defined in the API contract:
\begin{itemize}
    \item \texttt{GET /health}
    \item \texttt{POST /analyze} (multipart/form-data)
\end{itemize}

\subsection{Standard Response Envelope}

All API responses must be parsed using the standardized envelope.

\subsubsection{Success (Object)}
\begin{lstlisting}
{
  "success": true,
  "message": "Request completed successfully",
  "data": { },
  "meta": null
}
\end{lstlisting}

\subsubsection{Success (List + Pagination)}
\begin{lstlisting}
{
  "success": true,
  "message": "Request completed successfully",
  "data": [ ],
  "meta": {
    "pagination": {
      "page": 1,
      "limit": 10,
      "totalItems": 125,
      "totalPages": 13,
      "hasNextPage": true,
      "hasPreviousPage": false
    }
  }
}
\end{lstlisting}

\subsubsection{Error}
\begin{lstlisting}
{
  "success": false,
  "message": "Validation failed",
  "data": null,
  "meta": null,
  "error": {
    "code": "VALIDATION_ERROR",
    "message": "The uploaded file is missing the required column: gender."
  }
}
\end{lstlisting}

\subsection{Health Check Usage}

The frontend must include a lightweight health check mechanism:
\begin{itemize}
    \item Run \texttt{GET /health} on-demand (e.g., via a button) or as part of a diagnostics page/section.
    \item Display API availability status (online/offline) and the returned service label.
\end{itemize}

\subsection{Analyze Request (CSV Upload)}

\subsubsection{Request}

The frontend must submit analysis requests as \texttt{multipart/form-data} with:
\begin{itemize}
    \item \texttt{file}: the CSV file.
    \item \texttt{scoreType}: optional, defaults to \texttt{average}.
\end{itemize}

Supported \texttt{scoreType} values:
\begin{itemize}
    \item \texttt{average}, \texttt{math}, \texttt{reading}, \texttt{writing}
\end{itemize}

\subsubsection{Response}

The frontend must render the analysis response from:
\begin{itemize}
    \item \texttt{data.project}
    \item \texttt{data.analysis.scoreType}
    \item \texttt{data.analysis.groupVariable}
    \item \texttt{data.analysis.groups.male} and \texttt{data.analysis.groups.female}
    \item \texttt{data.analysis.comparison} fields
    \item \texttt{data.analysis.interpretation}
\end{itemize}

\section{User Experience Requirements}

\subsection{Primary Screens / Routes}

The frontend must include at minimum:
\begin{itemize}
    \item \textbf{Home / Upload}: Upload CSV + select scoreType + run analysis.
    \item \textbf{Results}: View results in a readable, structured layout (may be same page as upload).
    \item \textbf{About / Help}: Dataset format help, required columns, scoreType meaning, and API base URL.
\end{itemize}

\subsection{Upload Experience}

\begin{itemize}
    \item Allow selecting a local file via file picker and drag-and-drop (recommended).
    \item Show selected file name and size.
    \item Provide a clear ``Analyze'' call-to-action.
    \item Disable submit while a request is in progress.
    \item Show a loading state during analysis (spinner + message).
    \item Provide a ``Reset'' action to clear current results and selected file.
\end{itemize}

\subsection{Client-Side Validation (Preflight)}

Before sending the request, the frontend must validate:
\begin{itemize}
    \item A file is selected.
    \item File extension is \texttt{.csv} (basic validation).
    \item \texttt{scoreType} is one of the supported values.
\end{itemize}

If validation fails, show a user-friendly message and do not send the request.

\subsection{Error Handling}

The frontend must display errors consistently:
\begin{itemize}
    \item Use \texttt{message} as the primary user-facing summary.
    \item If present, show \texttt{error.code} and \texttt{error.message}.
    \item Provide an option to retry.
\end{itemize}

The UI must specifically handle common error codes from the API contract:
\begin{itemize}
    \item \texttt{NO\_FILE}
    \item \texttt{INVALID\_FILE\_TYPE}
    \item \texttt{MISSING\_COLUMNS}
    \item \texttt{INVALID\_DATA}
    \item \texttt{INSUFFICIENT\_GROUPS}
    \item \texttt{ANALYSIS\_FAILED}
\end{itemize}

\subsection{Results Presentation}

The results view must be structured into sections:
\begin{itemize}
    \item \textbf{Summary}: scoreType, groupVariable, sample sizes, significance.
    \item \textbf{Group Statistics}: male vs female metrics (mean, std dev, min, max).
    \item \textbf{Comparison / Inference}: meanDifference, confidence interval, pValue, tStatistic, degreesOfFreedom.
    \item \textbf{Interpretation}: display the provided interpretation text prominently.
\end{itemize}

Numerical formatting:
\begin{itemize}
    \item Use fixed decimal formatting for readability (e.g., 2 decimals), without changing the underlying values.
    \item Ensure p-values are shown clearly (e.g., \texttt{0.00005}).
\end{itemize}

\section{State Management and Caching}

\subsection{In-Memory State}

During a session, keep:
\begin{itemize}
    \item Selected file metadata (name, size).
    \item Selected scoreType.
    \item Latest analysis response.
    \item Loading and error states.
\end{itemize}

\subsection{Persistent Cache}

The frontend must cache the last analysis result in browser storage under:
\begin{lstlisting}
civara:last-analysis
\end{lstlisting}

Recommended cached object structure:
\begin{lstlisting}
{
  "uploadedAt": "2026-05-12T10:30:00.000Z",
  "fileName": "students.csv",
  "analysis": { }
}
\end{lstlisting}

Caching behavior:
\begin{itemize}
    \item Save on successful analysis response.
    \item Load on app start and restore the results view.
    \item Provide a UI action to clear the cache.
\end{itemize}

\section{Security and Privacy Requirements}

\begin{itemize}
    \item Do not log dataset contents to the console in production builds.
    \item Do not persist the raw CSV content in browser storage (store only API response + file name + timestamp).
    \item Use HTTPS in production deployments.
\end{itemize}

\section{Accessibility Requirements}

\begin{itemize}
    \item All form controls must have labels.
    \item Provide keyboard navigation for primary actions (upload, analyze, reset).
    \item Ensure color contrast for text and key UI elements.
\end{itemize}

\section{Performance Requirements}

\begin{itemize}
    \item Avoid blocking the UI thread during upload; show progress or loading indicators.
    \item Keep bundle size reasonable; avoid heavy charting libraries unless necessary.
\end{itemize}

\section{Acceptance Criteria}

The frontend is considered complete when:
\begin{itemize}
    \item A user can upload a CSV, select scoreType, and receive results via \texttt{POST /analyze}.
    \item Results render correctly using the standardized response envelope.
    \item Errors are displayed cleanly using \texttt{success}, \texttt{message}, and \texttt{error}.
    \item The latest successful analysis is cached and restored on refresh using \texttt{civara:last-analysis}.
    \item API base URL is configurable and defaults to production.
\end{itemize}

\end{document}

